\section{An independent signed-tabloid construction}
\label{l:tabloid-supplement}

The proof of Lemma~\ref{l:tabloid-embedding} in the main text uses the
row--column assignment from Lemma~\ref{l:small-index-character}.
Here we construct the signed tabloid directly from diagonal arms and
legs and prove its dimension bound from the hook formula. The same
part sizes then give an independent reciprocal-degree argument.
Throughout, $C_0=2{,}000{,}000$ as in
Theorem~\ref{l:tabloid-product}.

\begin{proof}[An arm--leg proof of Lemma~\ref{l:tabloid-embedding}]
Transposing $\lambda$ tensors its representation with sign and does
not change its dimension.  We may therefore orient the diagram so that
$\lambda_1\ge\lambda'_1$, where the prime denotes transpose, and undo
this orientation by a sign twist at the end.
Let $k$ be its diagonal length, and define
\[
 a_i=\lambda_i-i+1,\qquad b_i=\lambda'_i-i
 \quad(1\le i\le k).
\]
The $a_i$ count horizontal arms including their diagonal boxes; the
$b_i$ count vertical legs below those boxes.  They sum to $m$.
Use the sizes $a_1,b_1,\ldots,a_k,b_k$, omitting zero parts, with
trivial factors on the arms and sign factors on the legs.

To check occurrence, add arm $1$, leg $1$, arm $2$, leg $2$, and so
on.  Immediately before arm $i$, rows above $i$ have their final
lengths, while row $r\ge i$ has length $\min(i-1,\lambda_r)$.
Adding arm $i$ completes row $i$ and is a horizontal strip.
Adding leg $i$ then extends column $i$ down to height $\lambda'_i$
and is a vertical strip.  Each intermediate shape is a partition.
Repeated application of the horizontal and vertical Pieri rules
therefore shows that the resulting induced representation contains
$V_\lambda$.

It remains to compare its dimension with $D_\lambda$.  The hook-length
formula gives
\begin{equation}
 \frac{M}{D_\lambda}
 =\frac{\displaystyle\prod_{1\le i,j\le k}(a_i+b_j)}
 {\displaystyle\left(\prod_{i=1}^k a_i\right)
       \prod_{1\le i<j\le k}(a_i-a_j)(b_i-b_j)}.
 \label{l:tabloid-hook-ratio}
\end{equation}
Here is a direct verification of the identity.  The hooks in the
diagonal $k$ by $k$ square have lengths $a_i+b_j$.
For any partition with $k$ rows of lengths $r_i$, allowing zero rows,
the hook lengths in row $i$ are the integers from $1$ to $r_i+k-i$
with $r_i-r_j+j-i$ omitted for every $j>i$.  Indeed its hook at column
$t$ has length $r_i-t+1+\#\{j>i:r_j\ge t\}$; the gaps as $t$
increases are precisely those omitted values.  Apply this observation
to the diagram to the right of the square, whose row lengths are
$\lambda_i-k$, and to the transpose of the diagram below it, whose
row lengths are $\lambda'_i-k$.  Their hook products are, respectively,
\[
 \frac{\prod_i(a_i-1)!}{\prod_{i<j}(a_i-a_j)},
 \qquad
 \frac{\prod_i b_i!}{\prod_{i<j}(b_i-b_j)}.
\]
Division by $\prod_i a_i!b_i!$ proves
\eqref{l:tabloid-hook-ratio}.

When $k=1$, orientation gives $a_1\ge b_1+1$, so the ratio in
\eqref{l:tabloid-hook-ratio} is at most two.  A one-row diagram has
$M=D_\lambda=1$.  Every other such hook has $D_\lambda\ge2$, since
the entry $2$ can be placed either to the right of or below $1$ and
the resulting partial tableaux can be completed.  The claimed bound
follows in this case.

Assume now that $k\ge2$.  All logarithms in this proof are natural.
Arithmetic--geometric mean bounds the numerator in
\eqref{l:tabloid-hook-ratio} by $(m/k)^{k^2}$.  The two sequences
$a_i$ and $b_i$ are strictly decreasing integers, so their two
Vandermonde products are at least
$\prod_{j=1}^k((j-1)!)^2$.  For $1\le j\le k$,
\[
 \log((j-1)!)\ge (j-1)\log k-k.
\]
To see this, sum $\log(r/k)$ for $1\le r<j$, and bound that sum
below by the sum for $1\le r\le k$, which is at least $-k$ by
$k!\ge(k/e)^k$.  Consequently
\begin{equation}
 \log(M/D_\lambda)
 \le k^2\bigl(\log(m/k^2)+3\bigr).
 \label{l:tabloid-ratio-log}
\end{equation}

We need two lower bounds on $D_\lambda$ to pay for this ratio.
Every standard tableau of a contained diagram extends to $\lambda$
by adding the remaining boxes in a fixed order compatible with rows
and columns.  In particular the contained $k$ by $k$ square gives
\[
 D_\lambda\ge\frac{(k^2)!}{(2k)^{k^2}}
             \ge\left(\frac{k}{2e}\right)^{k^2}.
\]
It follows that
\begin{equation}
 \log D_\lambda\ge c_0 k^2\log k,
 \qquad c_0=10^{-5}.
 \label{l:tabloid-square-degree}
\end{equation}
For $k\ge64$ the displayed factorial estimate proves this directly.
For $2\le k<64$, the contained two-square gives $D_\lambda\ge2$,
and $10^{-5}63^2\log63<\log2$ suffices.

Orientation also gives $a_1\ge m/(2k)$, since every arm and leg has
size at most $a_1$.  The diagram contains $(a_1,k,\ldots,k)$ with
$k$ rows.  Fill its first $k$ top-row boxes with the smallest entries.
Then interleave the rest of the top row with any fixed standard order
of the $k(k-1)$ lower boxes.  Each interleaving is a standard tableau,
so
\[
 D_\lambda\ge
 \binom{a_1-k+k(k-1)}{k(k-1)}.
\]
If $m\ge16k^3$, put $r=m/k^3$, $u=a_1-k$ and $v=k(k-1)$.
Then $u\ge m/(4k)$ and $k^2/2\le v\le k^2$.  Using
$\binom{u+v}{v}\ge((u+v)/v)^v$ gives
\begin{equation}
 \log D_\lambda
 \ge\frac{k^2}{2}\log(r/4)
 \ge\frac{k^2}{4}\log r.
 \label{l:tabloid-long-arm-degree}
\end{equation}
If $m<16k^3$, then $\log(m/k^2)+3\le13\log k$, and
\eqref{l:tabloid-ratio-log}--\eqref{l:tabloid-square-degree} give
$\log M\le(1+13/c_0)\log D_\lambda$.
If $m\ge16k^3$, write
$\log(m/k^2)+3=\log(m/k^3)+\log k+3$ and use both degree bounds.
This gives
\[
 \log M\le
 \left(5+\frac{1+3/\log2}{c_0}\right)\log D_\lambda.
\]
Both coefficients are less than $2{,}000{,}000$, completing the proof.
\end{proof}

This arm--leg construction also gives an independent way to sum reciprocal
representation degrees. The same tabloid sizes therefore supply both
the embedding and the degree estimates used for counting small
restriction types and summing the final Fourier bounds.

\begin{lemma}[Degree sums from the tabloid embedding]
\label{l:tabloid-degree-sums}
For every fixed $s>0$,
\[
 \sum_{\substack{\lambda\vdash m\\D_\lambda>1}}
             D_\lambda^{-s}\longrightarrow0
 \qquad(m\longrightarrow\infty).
\]
In particular $\sum_{\lambda\vdash m}D_\lambda^{-1}\le3$ for all
sufficiently large $m$.
\end{lemma}

\begin{proof}
Let $p(u)$ count partitions of $u$, with $p(0)=1$.
For $u\ge1$, evaluating its generating product at $x=e^{-1/\sqrt u}$
gives
\[
 p(u)x^u\le\prod_{j\ge1}(1-x^j)^{-1},\qquad
 \log\prod_{j\ge1}(1-x^j)^{-1}
 =\sum_{i\ge1}\frac1{i(e^{i/\sqrt u}-1)}
 \le\sqrt u\sum_{i\ge1}i^{-2}.
\]
Thus $p(u)\le e^{3\sqrt u}$.
Orient $\lambda$ so that $\lambda_1\ge\lambda'_1$, and use the
preceding arm--leg construction with $a_1=\lambda_1$ and $j=m-a_1$.
For $1\le j\le m/2$, at most $2p(j)$ original
diagrams have this value: delete the oriented first row, and remember
whether a transpose was used.  The tabloid dimension satisfies
$M\ge\binom mj$, and that construction gives
\[
 D_\lambda^{-s}\le M^{-s/C_0}\le\binom mj^{-s/C_0}.
\]
For each fixed $j\ge1$ this tends to zero.  Since
$\binom mj\ge(m/j)^j\ge2^j$ on this range, the summable sequence
$2p(j)2^{-js/C_0}$ dominates the contributions.  They therefore sum
to $o(1)$.

If $j>m/2$, every arm or leg used in the tabloid has size at most
$a_1<m/2$.  Add some of these parts until their total $w$ first
reaches $m/4$.  Then $m/4\le w\le3m/4$ and
\[
 M\ge\binom mw\ge(m/w)^w\ge(4/3)^{m/4}.
\]
The total contribution of this range is at most
$p(m)(4/3)^{-ms/(4C_0)}=o(1)$.
The omitted $j=0$ diagrams are the row and column, whose
representations are trivial and sign.  They are the only
one-dimensional representations for $m\ge2$, and their contribution
to the inverse-first-degree sum is two.
\end{proof}
